% ============================================================================
%  SECCIÓN 1.7: DELIMITACIÓN
% ============================================================================

\section{Delimitaci\'on}

La delimitaci\'on, entendida como la acci\'on y efecto de delimitar, hace
referencia a determinar las caracter\'isticas esenciales del contexto y del
objeto de estudio. En el proyecto de grado, la delimitaci\'on del contexto y del
objeto mismo de estudio es primordial, pues permite situar el problema y las
razones del mismo en el contexto especificado \citep{arias2012}.

\subsection{L\'imite temporal}

El proyecto se desarrollar\'a durante la gesti\'on acad\'emica correspondiente a
la asignatura de Aplicaciones M\'oviles I, en el marco del aprendizaje basado en
proyectos (ABP), con una duraci\'on estimada de seis meses, comprendidos entre
agosto de 2026 y enero de 2027. Las actividades se distribuyen en un
cronograma de 16 semanas organizadas en cinco fases: an\'alisis, dise\~no,
desarrollo, pruebas y cierre, tal como se detalla en el anexo correspondiente.

\subsection{L\'imite geogr\'afico}

El \'ambito geogr\'afico del proyecto es el territorio nacional, con \'enfasis en
los escenarios deportivos de las principales ciudades del pa\'is. La
informaci\'on se incorporar\'a de forma progresiva, comenzando por un piloto de
validaci\'on en las ciudades de La Paz, Cochabamba y Santa Cruz de la Sierra,
donde se concentra la mayor oferta de infraestructura deportiva del pa\'is.

\subsection{Otros l\'imites del proyecto}

El proyecto se limita a la visualizaci\'on interactiva y a la difusi\'on de
informaci\'on sobre los escenarios deportivos. Quedan fuera del alcance los
procesos de reserva, alquiler y pago en l\'inea, as\'i como la gesti\'on
administrativa interna de los escenarios, que podr\'an ser considerados en
avances futuros.
